% Target: rmp-workbench-iii-historical-composite
% Formula: Historical composite of the four rotated boundary actions.
% Ink: Kernel flat frame; the four-station boundary loop is one closed operator
%   string through the lambda corner, and each sweeping arc is an operator
%   string routed outside the loop.
% Boundary: Resolved by the public environment; archival boundary contract is explicit.
% Source: author source ImagesReview Section III/ImagesReview Section III.tex lines
%   170-197; archival output ImagesReview Section III.pdf
% Capabilities: kernel, path-string, typed-ports
\tenkzkernel
\tndeclare{species}{operator}{hue=source:red}
\[
  % LHS: the two pendant actions feed INTO the loop and the loop's own
  % half-edges leave outward; the boundary string sweeps around the top-left,
  % outside the loop.
  \begin{tenkz}[rows={wire,wire,wire,wire}, cols=4, bonds=none, boundary=none]
    \tn[at=(1,3), skin=box, name=top, ports={90:virtual, 270:virtual}]{}
    \tn[at=(2,2), skin=ring, name=lam, label pos=nw]{\lambda}
    \tn[at=(2,3), skin=box, name=n, ports={90:virtual}]{}
    \tn[at=(3,1), skin=box, name=left, ports={0:virtual, 180:virtual}]{}
    \tn[at=(3,2), skin=box, name=w, ports={180:virtual}]{}
    \tn[at=(3,4), skin=box, name=e, ports={0:virtual}]{}
    \tn[at=(4,3), skin=box, name=s, ports={270:virtual}]{}
    % the closed boundary loop, chained through the four legged atoms and
    % three invisible corner junctions (via= retired, #6897)
    \tn[at=(2,4), skin=none, name=jring1]{}
    \tn[at=(4,4), skin=none, name=jring2]{}
    \tn[at=(4,2), skin=none, name=jring3]{}
    \tnwire[kind=string, species=operator, name=ring]{lam}{n}
    \tnwire[kind=string, species=operator, name=r2]{n}{jring1}
    \tnwire[kind=string, species=operator, name=r3,
      cross={under at crossing of r2 and r3}]{jring1}{e}
    \tnwire[kind=string, species=operator, name=r4]{e}{jring2}
    \tnwire[kind=string, species=operator, name=r5,
      cross={under at crossing of r4 and r5}]{jring2}{s}
    \tnwire[kind=string, species=operator, name=r6]{s}{jring3}
    \tnwire[kind=string, species=operator, name=r7,
      cross={under at crossing of r6 and r7}]{jring3}{w}
    \tnwire[kind=string, species=operator, name=ringclose,
      cross={under at crossing of ring and ringclose}]{w}{lam}
    \tnwire[dir=to]{top.270}{n.90} \tnwire{top.90}{open n}
    \tnwire[dir=to]{left.0}{w.180} \tnwire{left.180}{open w}
    \tnwire[dir=to]{s.270}{open s} \tnwire[dir=to]{e.0}{open e}
    % the sweeping arc, chained through one invisible junction near lam
    \tn[at=1 nw of lam, skin=none, name=jarm1]{}
    \tnwire[kind=string, species=operator, name=arm]{left}{jarm1}
    \tnwire[kind=string, species=operator, name=arm2,
      cross={under at crossing of arm and arm2}]{jarm1}{top}
  \end{tenkz}
  \;=\;
  % RHS: the loop's half-edges arrive from outward and the two pendant
  % actions drain OUT of the loop; the boundary string sweeps around the
  % bottom-right.
  \begin{tenkz}[rows={wire,wire,wire,wire}, cols=4, bonds=none, boundary=none]
    \tn[at=(1,1), skin=ring, name=lam, label pos=nw]{\lambda}
    \tn[at=(1,2), skin=box, name=n, ports={90:virtual}]{}
    \tn[at=(2,1), skin=box, name=w, ports={180:virtual}]{}
    \tn[at=(2,3), skin=box, name=e, ports={0:virtual}]{}
    \tn[at=(2,4), skin=box, name=right, ports={0:virtual, 180:virtual}]{}
    \tn[at=(3,2), skin=box, name=s, ports={270:virtual}]{}
    \tn[at=(4,2), skin=box, name=bottom, ports={90:virtual, 270:virtual}]{}
    \tn[at=(1,3), skin=none, name=jring1]{}
    \tn[at=(3,3), skin=none, name=jring2]{}
    \tn[at=(3,1), skin=none, name=jring3]{}
    \tnwire[kind=string, species=operator, name=ring]{lam}{n}
    \tnwire[kind=string, species=operator, name=r2]{n}{jring1}
    \tnwire[kind=string, species=operator, name=r3,
      cross={under at crossing of r2 and r3}]{jring1}{e}
    \tnwire[kind=string, species=operator, name=r4]{e}{jring2}
    \tnwire[kind=string, species=operator, name=r5,
      cross={under at crossing of r4 and r5}]{jring2}{s}
    \tnwire[kind=string, species=operator, name=r6]{s}{jring3}
    \tnwire[kind=string, species=operator, name=r7,
      cross={under at crossing of r6 and r7}]{jring3}{w}
    \tnwire[kind=string, species=operator, name=ringclose,
      cross={under at crossing of ring and ringclose}]{w}{lam}
    \tnwire[dir=to]{open n}{n.90} \tnwire[dir=to]{open w}{w.180}
    \tnwire[dir=to]{s.270}{bottom.90} \tnwire{bottom.270}{open s}
    \tnwire[dir=to]{e.0}{right.180} \tnwire{right.0}{open e}
    \tn[at=1 se of (3,3), skin=none, name=jarm1]{}
    \tnwire[kind=string, species=operator, name=arm]{right}{jarm1}
    \tnwire[kind=string, species=operator, name=arm2,
      cross={under at crossing of arm and arm2}]{jarm1}{bottom}
  \end{tenkz}
\]
